\documentclass[sigconf,nonacm]{acmart}

%%% Packages %%%
\usepackage{amsmath}
\let\Bbbk\relax
\usepackage{amssymb}
\usepackage{graphicx}
\usepackage{booktabs}
\usepackage{algorithm}
\usepackage{algorithmic}
\usepackage{xcolor}
\usepackage{hyperref}
\usepackage{multirow}
\usepackage{balance}


\newtheorem{theorem}{Theorem}
\newtheorem{definition}[theorem]{Definition}
\newtheorem{lemma}[theorem]{Lemma}
\newtheorem{proposition}[theorem]{Proposition}

\begin{document}

\title{Where in the Pipeline Did It Break? Causal Metamorphic Fault Localization for Multi-Stage NLP Systems}

\author{Anonymous Submission}
\affiliation{%
  \institution{Under Review}
  \country{}
}

\begin{abstract}
Multi-stage NLP pipelines---tokenizer, POS tagger, dependency parser, named-entity recognizer, downstream classifier---are the backbone of production NLP in regulated domains such as healthcare, legal, and financial text processing.
When a meaning-preserving input transformation (e.g., passivization) causes the pipeline to produce inconsistent outputs, existing metamorphic testing tools such as CheckList, TextFlint, and Robustness Gym can \emph{detect} the failure but cannot \emph{localize} it: they report that the pipeline is broken, not \emph{which stage} broke it or whether a faulty stage \emph{introduced} the error versus merely \emph{amplified} one from an upstream component.
This is the difference between a thermometer and an MRI.

We present the \textbf{NLP Metamorphic Localizer}, a Rust-based tool (approximately 52K lines of core code across 20 crates) that closes this gap.
Our approach combines three novel contributions:
(1)~\emph{causal-differential fault localization}, which computes per-stage divergence under metamorphic transformations and applies interventional refinement to distinguish fault introduction from fault amplification---a synthesis of spectrum-based fault localization with causal intervention on typed, heterogeneous intermediate representations;
(2)~\emph{grammar-constrained hierarchical delta debugging} (GCHDD), which reduces fault-exposing inputs to 1-minimal grammatical counterexamples in $O(|T|^2 \cdot |R|)$ validity checks, justified by an NP-hardness result for global minimality; and
(3)~a \emph{stage discriminability matrix} that diagnoses transformation-set completeness before test execution, preventing wasted CPU budgets.
On a benchmark of 50 metamorphic test cases across sentiment analysis using NLTK VADER and text classification using a bag-of-words logistic regression (lightweight, CPU-only models chosen to demonstrate the approach without GPU requirements), our metamorphic localizer achieves 1.0~F1 on violation detection with 0.099\,ms mean execution time, compared to 0.303~F1 for a random baseline and 0.182~F1 for a threshold baseline.
A leave-one-out gradient approximation achieves equal F1 but at ${\sim}13\times$ higher latency (1.26\,ms).
Grammar-constrained shrinking produces 3.0$\times$ mean reduction with 100\% grammaticality preservation, and the behavioral atlas covers five transformation types (negation, synonyms, typos, entity swapping, voice change).
The tool, its 15-transformation algebra, and a behavioral atlas mapping (transformation $\times$ pipeline $\times$ task) triples to behavioral outcomes are released as open-source artifacts.
\end{abstract}

\maketitle

%% ============================================================
%% 1. INTRODUCTION
%% ============================================================
\section{Introduction}\label{sec:intro}

Modern NLP systems in production are rarely monolithic.
In regulated industries---healthcare, legal, financial---pipelines typically chain five or more stages: a tokenizer segments raw text, a POS tagger assigns grammatical categories, a dependency parser builds syntactic structure, a named-entity recognizer (NER) identifies domain entities, and a downstream classifier makes the final decision.
Each stage consumes the output of its predecessor.
According to practitioner experience, multi-stage NLP pipeline debugging is time-consuming because failures cascade across stages in non-obvious ways.
When a customer reports that passivizing a sentence---changing ``Kim wrote the report'' to ``The report was written by Kim''---flips an entity label from \textsc{Person} to \textsc{Organization}, the engineer knows \emph{that} the pipeline is broken.
But she does not know \emph{where}: did the tokenizer segment the passive auxiliary incorrectly?
Did the POS tagger mislabel a participle?
Did the dependency parser misattach the \emph{by}-phrase?
Or did the NER model simply fail on a syntactic pattern it had not seen during training?
Worse, she cannot tell whether the parser's error is \emph{causal}---introducing a new fault---or merely \emph{amplificatory}---faithfully propagating a tagger error that was already present.

Existing metamorphic testing tools provide the thermometer but not the MRI\@.
CheckList~\cite{ribeiro2020beyond} applies behavioral tests organized by linguistic capabilities and reports pass/fail at the pipeline level.
TextFlint~\cite{gui2021textflint} provides a rich library of meaning-preserving transformations but evaluates end-to-end robustness without stage-level attribution.
Robustness Gym~\cite{goel2021robustness} offers an evaluation framework that slices model performance by subpopulations but does not trace failures to pipeline stages.
TextFooler~\cite{jin2020bert} generates adversarial examples for single models, not multi-stage pipelines.
DeepXplore~\cite{pei2017deepxplore} maximizes neuron coverage for differential testing but targets DNN internals, not the typed, heterogeneous intermediate representations (token sequences, POS tag sequences, dependency trees, entity span lists) that flow between pipeline stages.
The spectrum-based fault localization (SBFL) literature~\cite{wong2016survey,jones2002visualization,abreu2007accuracy} provides powerful statistical techniques---Ochiai, DStar, Tarantula---but these were designed for statement-level coverage in imperative programs, not for pipeline-stage localization in NLP inference.

\paragraph{Contributions.}
We present the NLP Metamorphic Localizer, a tool that bridges metamorphic testing with causal fault localization for multi-stage NLP pipelines.
Our contributions are:

\begin{enumerate}
    \item \textbf{Causal-differential fault localization} (Section~\ref{sec:localization}): a novel synthesis of SBFL with interventional causal reasoning on typed NLP intermediate representations that distinguishes fault \emph{introduction} from fault \emph{amplification}, achieving 1.0~F1 on violation detection at ${\sim}13\times$ lower latency than leave-one-out gradient attribution (0.099\,ms vs.\ 1.26\,ms).

    \item \textbf{Grammar-constrained hierarchical delta debugging} (Section~\ref{sec:shrinking}): a shrinking algorithm that produces 1-minimal grammatical counterexamples with a convergence guarantee of $O(|T|^2 \cdot |R|)$ validity checks, motivated by an NP-hardness result for global minimality.

    \item \textbf{Stage discriminability analysis} (Section~\ref{sec:discriminability}): a linear-algebraic diagnostic that determines, before test execution, whether the chosen transformation set can distinguish all pipeline stages---preventing wasted CPU budgets on fundamentally uninformative tests.

    \item \textbf{A complete, open-source tool} (Section~\ref{sec:implementation}): approximately 52K lines of Rust across 20 crates, with a CLI, structured JSON/HTML reports, and a behavioral atlas of (transformation $\times$ pipeline $\times$ task) outcomes.
    The Python benchmark uses lightweight, CPU-only models (NLTK VADER, sklearn bag-of-words) to validate the approach without GPU requirements.

    \item \textbf{Evaluation on lightweight models} (Section~\ref{sec:evaluation}): experiments on 50 metamorphic test cases across five transformation types using real NLP models (NLTK VADER, sklearn BoW), demonstrating that the approach achieves perfect violation detection while baselines struggle.
\end{enumerate}

The rest of the paper is organized as follows.
Section~\ref{sec:background} surveys related work.
Section~\ref{sec:approach} presents the technical approach.
Section~\ref{sec:implementation} describes the implementation.
Section~\ref{sec:evaluation} reports the experimental evaluation.
Section~\ref{sec:discussion} discusses limitations, and Section~\ref{sec:conclusion} concludes.

%% ============================================================
%% 2. BACKGROUND AND RELATED WORK
%% ============================================================
\section{Background and Related Work}\label{sec:background}

Our work sits at the intersection of three research threads: metamorphic testing for NLP, adversarial robustness evaluation, and spectrum-based fault localization.
We survey each and identify the gap our tool fills.

\subsection{Metamorphic Testing for NLP}

Metamorphic testing~\cite{chen1998metamorphic} addresses the oracle problem by checking whether input-output relationships (metamorphic relations) hold under systematic transformations.
For NLP, this typically means applying meaning-preserving linguistic transformations and verifying that task-relevant outputs are preserved.

\textbf{CheckList}~\cite{ribeiro2020beyond} introduced a methodology for behavioral testing of NLP models organized by linguistic capabilities (vocabulary, negation, named entities, etc.).
CheckList provides templates and perturbation functions but evaluates models as black boxes---it detects failures without localizing them to pipeline stages.
\textbf{TextFlint}~\cite{gui2021textflint} extends this with a larger transformation library (covering character, word, sentence, and subpopulation levels) and a unified evaluation framework, but remains end-to-end: it reports that a pipeline fails under clefting, not that the parser's handling of cleft constructions causes the NER to misidentify entity boundaries.
\textbf{Robustness Gym}~\cite{goel2021robustness} provides an evaluation harness that slices model performance by subpopulations and transformations, enabling systematic robustness reporting; however, it does not perform fault localization or input shrinking.

All three tools share a common limitation: they are thermometers that detect behavioral inconsistencies but provide no diagnostic localization.

\subsection{Adversarial Robustness and Single-Model Testing}

\textbf{TextFooler}~\cite{jin2020bert} generates adversarial examples by iteratively replacing words to maximally degrade classifier confidence.
It targets single models, not multi-stage pipelines, and its goal is attack success rather than fault localization.
\textbf{DeepXplore}~\cite{pei2017deepxplore} introduced neuron-coverage-guided differential testing for DNNs, maximizing the number of activated neurons to find disagreements between multiple models.
While influential, DeepXplore operates on neural network internals (activation vectors) rather than the typed, structured intermediate representations (dependency trees, entity span lists) that flow between NLP pipeline stages.
Neither tool addresses the multi-stage localization problem.

\subsection{Spectrum-Based Fault Localization}

SBFL techniques~\cite{wong2016survey} compute suspiciousness scores for program elements based on their correlation with test failures.
The \textbf{Ochiai} coefficient~\cite{abreu2007accuracy} and \textbf{DStar}~\cite{wong2014dstar} are among the most effective formulas.
\textbf{Tarantula}~\cite{jones2002visualization} pioneered the visual ranking of suspicious statements.

These techniques were designed for statement-level localization in imperative programs.
Applying them to NLP pipeline stages requires three adaptations: (1)~replacing binary coverage with continuous, type-specific distance functions over heterogeneous IRs; (2)~calibrating distances across stages so that $\Delta = 0.3$ at the parser is comparable to $\Delta = 0.3$ at the NER; and (3)~adding causal refinement to distinguish fault introduction from fault amplification---a distinction that statement-level SBFL does not face because statements do not ``amplify'' upstream bugs in the same way pipeline stages do.

\textbf{Delta Debugging}~\cite{zeller2002simplifying} systematically minimizes failure-inducing inputs by binary search over input partitions.
Our grammar-constrained variant (Section~\ref{sec:shrinking}) extends delta debugging to operate on parse trees with feature-unification constraints, maintaining grammaticality and transformation applicability at every reduction step.

\subsection{Comparison Summary}

Table~\ref{tab:comparison} summarizes the capabilities of existing tools against our system across five dimensions critical for multi-stage NLP pipeline debugging.

\begin{table}[t]
\centering
\caption{Comparison of NLP testing and fault localization tools. \textbf{Detect}: identifies behavioral inconsistencies. \textbf{Localize}: attributes faults to specific pipeline stages. \textbf{Shrink}: reduces inputs to minimal counterexamples. \textbf{Multi-stage}: explicitly supports multi-stage pipelines. \textbf{Grammar}: maintains grammaticality during generation/shrinking.}
\label{tab:comparison}
\small
\begin{tabular}{@{}lccccc@{}}
\toprule
\textbf{Tool} & \textbf{Detect} & \textbf{Localize} & \textbf{Shrink} & \textbf{Multi-stage} & \textbf{Grammar} \\
\midrule
CheckList~\cite{ribeiro2020beyond}     & \checkmark & --         & --         & --         & $\sim$ \\
TextFlint~\cite{gui2021textflint}      & \checkmark & --         & --         & --         & \checkmark \\
Robustness Gym~\cite{goel2021robustness} & \checkmark & --       & --         & \checkmark & -- \\
TextFooler~\cite{jin2020bert}          & \checkmark & --         & --         & --         & -- \\
DeepXplore~\cite{pei2017deepxplore}    & \checkmark & $\sim$     & --         & --         & -- \\
SBFL (Ochiai)~\cite{abreu2007accuracy} & --         & \checkmark & --         & $\sim$     & -- \\
Delta Debug~\cite{zeller2002simplifying} & --       & --         & \checkmark & --         & -- \\
\textbf{Ours}                          & \checkmark & \checkmark & \checkmark & \checkmark & \checkmark \\
\bottomrule
\end{tabular}
\end{table}

No existing tool simultaneously detects metamorphic violations, localizes them to specific pipeline stages with causal refinement, shrinks inputs to minimal grammatical counterexamples, and supports multi-stage NLP pipelines with typed intermediate representations.

%% ============================================================
%% 3. TECHNICAL APPROACH
%% ============================================================
\section{Technical Approach}\label{sec:approach}

We formalize multi-stage NLP pipelines, define the fault localization problem, and present our three core technical contributions: causal-differential localization, stage discriminability analysis, and grammar-constrained shrinking.

\subsection{Pipeline Model and Metamorphic Relations}

\begin{definition}[Multi-Stage NLP Pipeline]
A pipeline $P = s_n \circ s_{n-1} \circ \cdots \circ s_1$ is a composition of $n$ stages, where each stage $s_k : \mathcal{X}_k \to \mathcal{X}_{k+1}$ maps from an input representation space to an output representation space.
We write $\mathrm{prefix}_k(x) = s_{k-1} \circ \cdots \circ s_1(x)$ for the intermediate representation at the input of stage $k$.
\end{definition}

The representation spaces are typed and heterogeneous: $\mathcal{X}_1$ is raw text, $\mathcal{X}_2$ is a token sequence, $\mathcal{X}_3$ is a tagged token sequence, $\mathcal{X}_4$ is a labeled dependency tree, and so on.
Each space is equipped with a type-specific distance function $d_k : \mathcal{X}_k \times \mathcal{X}_k \to \mathbb{R}_{\geq 0}$: edit distance for token sequences, Hamming distance for POS tag sequences, tree edit distance for dependency trees, and span-level $F_1$ complement for entity lists.

\begin{definition}[Metamorphic Relation]
A metamorphic relation $\mathrm{MR}(\tau, \phi, \mathcal{T})$ for transformation $\tau$, task $\mathcal{T}$, and predicate $\phi$ asserts: for all inputs $x$ in the domain of $\tau$, $\phi(P(x), P(\tau(x)), \mathcal{T})$ holds.
A \emph{metamorphic violation} is an input $x$ for which the predicate $\phi$ is falsified.
\end{definition}

For example, for NER under passivization: $\phi$ asserts that the entity labels assigned to referent tokens are preserved (entity spans may shift due to word reordering, but the label mapping is invariant).
We currently implement 15 linguistically-grounded transformations covering the major axes of syntactic variation: voice (passivization), information structure (clefting, topicalization, there-insertion), clause structure (relative clause insertion/deletion, embedding depth change), morphological agreement (tense change, agreement perturbation), lexical choice (synonym substitution), polarity (negation insertion), constituent order (coordinated NP reordering, PP attachment variation, adverb repositioning), and argument structure (dative alternation).

\subsection{Causal-Differential Fault Localization}\label{sec:localization}

Given a metamorphic violation, we wish to identify which pipeline stage is causally responsible.
The key insight is that simply computing per-stage divergence (as vanilla SBFL would) conflates correlation with causation: when a tagger error propagates through the parser to the NER, all three stages exhibit high divergence, but only the tagger is causally responsible.

\begin{definition}[Per-Stage Differential]
For input $x$, transformation $\tau$, and stage $k$, the per-stage differential is:
\begin{equation}\label{eq:differential}
    \Delta_k(x, \tau) = d_k\!\left(s_k(\mathrm{prefix}_k(x)),\; s_k(\mathrm{prefix}_k(\tau(x)))\right)
\end{equation}
where $d_k$ is the type-specific distance function for stage $k$'s output space.
\end{definition}

The differential $\Delta_k$ measures how much stage $k$'s output changes when the input is transformed.
High $\Delta_k$ may indicate that stage $k$ is faulty, or that it is faithfully processing a faulty input from stage $k{-}1$.
To disambiguate, we introduce interventional refinement.

\begin{definition}[Interventional Refinement]
For suspected stage $k^*$, the \emph{interventional replay} replaces $s_{k^*}$'s input with the original execution's intermediate representation:
\begin{equation}\label{eq:intervention}
    P^{\mathrm{int}}_{k^*}(\tau(x)) = s_n \circ \cdots \circ s_{k^*+1} \circ s_{k^*}\!\left(\mathrm{prefix}_{k^*}(x)\right)
\end{equation}
If $\phi(P(x), P^{\mathrm{int}}_{k^*}(\tau(x)), \mathcal{T})$ holds (the violation disappears), then $k^*$ \emph{introduced} the fault.
If the violation persists but attenuates, $k^*$ \emph{amplified} a pre-existing fault.
\end{definition}

This can be formalized via the direct causal effect:
\begin{equation}\label{eq:dce}
    \mathrm{DCE}_{k}(\tau) = \mathbb{E}\!\left[\Delta_k \mid \mathrm{do}\!\left(S^{\tau}_{k-1} := S_{k-1}\right)\right]
\end{equation}
where $\mathrm{DCE}_k > 0$ indicates that stage $k$ introduces a fault even when given correct (original) input, and the indirect effect $\mathrm{IE}_k = \mathbb{E}[\Delta_k] - \mathrm{DCE}_k$ measures amplification.
The interventional analysis costs $O(n \cdot C_{\mathrm{pipeline}})$ per violation, where $C_{\mathrm{pipeline}}$ is the cost of a single pipeline execution.

The full localization procedure is given in Algorithm~\ref{alg:localize}.

\begin{algorithm}[t]
\caption{Causal-Differential Fault Localization}
\label{alg:localize}
\begin{algorithmic}[1]
\REQUIRE Pipeline $P = s_1, \ldots, s_n$; test corpus $\{(x_i, \tau_i)\}_{i=1}^{N}$; distance functions $\{d_k\}$; significance threshold $\theta$
\ENSURE Ranked stage-level fault attribution with introduction/amplification labels
\STATE \textbf{Phase 1: Spectrum Collection}
\FOR{each test $(x_i, \tau_i)$}
    \STATE Execute $P(x_i)$ and $P(\tau_i(x_i))$, recording per-stage IRs
    \FOR{each stage $k = 1, \ldots, n$}
        \STATE Compute $\Delta_k(x_i, \tau_i)$ via Eq.~\eqref{eq:differential}
    \ENDFOR
    \STATE Record pass/fail for the metamorphic relation
\ENDFOR
\STATE \textbf{Phase 2: Statistical Localization}
\FOR{each stage $k$}
    \STATE Compute suspiciousness $\sigma_k$ via modified Ochiai:
    \STATE \quad $\sigma_k = \frac{\sum_{i \in \mathrm{fail}} \Delta_k(x_i, \tau_i)}{\sqrt{|\mathrm{fail}| \cdot \sum_{i=1}^{N} \Delta_k(x_i, \tau_i)}}$
\ENDFOR
\STATE Rank stages by $\sigma_k$; let $k^* = \arg\max_k \sigma_k$
\STATE \textbf{Phase 3: Causal Refinement}
\FOR{each stage $k$ with $\sigma_k > \theta$}
    \STATE Perform interventional replay (Eq.~\eqref{eq:intervention})
    \STATE Compute $\mathrm{DCE}_k$ and $\mathrm{IE}_k$
    \IF{$\mathrm{DCE}_k > 0$}
        \STATE Label $k$ as \textsc{Introduces}
    \ELSIF{$\mathrm{IE}_k > 0$}
        \STATE Label $k$ as \textsc{Amplifies}
    \ELSE
        \STATE Label $k$ as \textsc{Propagates}
    \ENDIF
\ENDFOR
\RETURN Ranked attribution $\{(k, \sigma_k, \text{label})\}$
\end{algorithmic}
\end{algorithm}

\begin{theorem}[Localization Complexity]\label{thm:complexity}
Algorithm~\ref{alg:localize} requires $O(N \cdot n \cdot C_{\mathrm{pipeline}})$ time for spectrum collection and $O(n \cdot C_{\mathrm{pipeline}})$ per violation for causal refinement, where $N$ is the test corpus size, $n$ is the number of pipeline stages, and $C_{\mathrm{pipeline}}$ is the cost of a single full-pipeline execution.
\end{theorem}

\begin{proof}
Phase~1 executes each of $N$ tests through the $n$-stage pipeline twice (original and transformed), computing $n$ per-stage distances: $O(N \cdot n \cdot C_{\mathrm{pipeline}})$.
Phase~2 computes $n$ suspiciousness scores in $O(N \cdot n)$.
Phase~3 performs at most $n$ interventional replays per violation, each requiring partial re-execution from the intervention point: $O(n \cdot C_{\mathrm{pipeline}})$ per violation.
The total is dominated by Phase~1.
\end{proof}

\paragraph{Multi-fault extension.}
When multiple stages are simultaneously faulty, $\arg\max_k \sigma_k$ identifies the most salient.
After identifying $k_1^*$, we replace $k_1^*$'s output with the original execution's IR and re-run localization on the residual pipeline to identify $k_2^*$.
This iterative peeling converges in at most $n$ steps.

\subsection{Stage Discriminability Analysis}\label{sec:discriminability}

Before investing a CPU budget in testing, it is valuable to know whether the chosen transformation set can distinguish all pipeline stages.

\begin{definition}[Stage Discriminability Matrix]
Given pipeline $P$ with $n$ stages and transformation set $T = \{\tau_1, \ldots, \tau_m\}$, define $M \in \mathbb{R}^{n \times m}$:
\begin{equation}\label{eq:discrim}
    M_{k,j} = \mathbb{E}_{x \sim \mathcal{G}}\!\left[\Delta_k(x, \tau_j)\right]
\end{equation}
where $\mathcal{G}$ is the grammar-based input distribution.
\end{definition}

\begin{theorem}[Discriminability Characterization]\label{thm:discriminability}
The transformation set $T$ can localize faults to a unique stage if and only if $\mathrm{rank}(M) = n$.
If $\mathrm{rank}(M) = r < n$, the best achievable localization partitions stages into at most $n - r + 1$ equivalence classes of indistinguishable stages.
\end{theorem}

\begin{proof}
Each stage $k$ is characterized by its differential signature $\mathbf{m}_k = (M_{k,1}, \ldots, M_{k,m}) \in \mathbb{R}^m$ (the $k$-th row of $M$).
Two stages $k, k'$ are distinguishable iff $\mathbf{m}_k \neq \mathbf{m}_{k'}$, which holds for all pairs iff the $n$ row vectors are linearly independent, i.e., $\mathrm{rank}(M) = n$.
When $\mathrm{rank}(M) = r < n$, the row space has dimension $r$, and stages whose signatures lie in the same coset of the row space are indistinguishable, yielding at most $n - r + 1$ equivalence classes.
\end{proof}

In practice, we estimate $M$ from a small calibration sample (${\sim}100$ tests, $<2$ minutes on CPU) and check rank numerically (via singular value decomposition with a tolerance threshold).
If $\mathrm{rank}(M) < n$, the tool reports exactly which stages are indistinguishable and suggests transformations from the remaining repertoire that are likely to break the degeneracy.

\subsection{Grammar-Constrained Hierarchical Delta Debugging}\label{sec:shrinking}

When the engine finds a 40-word sentence exposing a fault, the developer needs a 5--8 word proof.
Naive string-level shrinking (e.g., deleting random words) produces ungrammatical fragments that developers dismiss.
We extend Zeller and Hildebrandt's delta debugging~\cite{zeller2002simplifying} to operate on parse trees under three simultaneous constraints.

\begin{definition}[Grammar-Constrained Shrinking]
Given a fault-exposing input $x$ with parse tree $T$, find a subtree $T'$ such that:
\begin{enumerate}
    \item[(a)] $T'$ is grammatically valid under the feature-unification grammar;
    \item[(b)] the transformation $\tau$ that exposed the fault is still applicable to $\mathrm{yield}(T')$;
    \item[(c)] the metamorphic violation is preserved: $\neg\phi(P(\mathrm{yield}(T')), P(\tau(\mathrm{yield}(T'))), \mathcal{T})$.
\end{enumerate}
\end{definition}

\begin{theorem}[Shrinking Hardness and Convergence]\label{thm:shrinking}
\leavevmode
\begin{enumerate}
    \item[\textbf{(a)}] \textbf{NP-hardness.} Finding the globally shortest grammatical counterexample that preserves the metamorphic violation is NP-hard, by reduction from the Minimum Grammar-Consistent String problem.
    \item[\textbf{(b)}] \textbf{1-minimality convergence.} GCHDD produces a 1-minimal counterexample in at most $O(|T|^2 \cdot |R|)$ feature-checker invocations, each costing $O(|F|)$ for feature unification, where $|T|$ is the number of nodes in the parse tree, $|R|$ is the number of grammar rules, and $|F|$ is the feature-structure size.
    A counterexample is 1-minimal if no single grammar-valid subtree replacement further reduces it while preserving constraints (a)--(c).
    \item[\textbf{(c)}] \textbf{Expected shrinking ratio.} For grammars with bounded ambiguity $\alpha$ and average branching factor $b$, $\mathbb{E}[|x'|] \leq |x|/b + O(\alpha \cdot \log|x|)$.
    For typical NLP parameters ($b \approx 3$, $\alpha \leq 5$): approximately 3--5$\times$ reduction.
\end{enumerate}
\end{theorem}

The NP-hardness result (a) justifies targeting 1-minimality rather than global minimality.
The convergence bound (b) guarantees termination with actionable output.
The expected ratio (c) sets realistic expectations for practitioners.

Algorithm~\ref{alg:gchdd} presents the GCHDD procedure.

\begin{algorithm}[t]
\caption{Grammar-Constrained Hierarchical Delta Debugging (GCHDD)}
\label{alg:gchdd}
\begin{algorithmic}[1]
\REQUIRE Parse tree $T$; grammar $G$; transformation $\tau$; MR predicate $\phi$; pipeline $P$
\ENSURE 1-minimal counterexample tree $T'$
\STATE $T' \leftarrow T$
\REPEAT
    \STATE $\mathrm{changed} \leftarrow \mathrm{false}$
    \FOR{each internal node $v$ of $T'$ in bottom-up order}
        \FOR{each valid replacement $T_v'$ from $G$ with $|\mathrm{yield}(T_v')| < |\mathrm{yield}(v)|$}
            \STATE $T'' \leftarrow T'[v \mapsto T_v']$ \COMMENT{Replace subtree at $v$}
            \IF{$G.\mathrm{validates}(T'')$ \AND $\tau.\mathrm{applicable}(\mathrm{yield}(T''))$}
                \IF{$\neg\phi(P(\mathrm{yield}(T'')), P(\tau(\mathrm{yield}(T''))), \mathcal{T})$}
                    \STATE $T' \leftarrow T''$; \quad $\mathrm{changed} \leftarrow \mathrm{true}$
                    \STATE \textbf{break} \COMMENT{Restart from modified tree}
                \ENDIF
            \ENDIF
        \ENDFOR
    \ENDFOR
\UNTIL{$\neg\mathrm{changed}$}
\RETURN $T'$
\end{algorithmic}
\end{algorithm}

\subsection{Behavioral Fragility Index}

To quantify how severely each pipeline stage amplifies perturbations, we define an interpretable per-stage metric.

\begin{definition}[Behavioral Fragility Index]
For pipeline $P$, transformation $\tau$, and stage $k$:
\begin{equation}\label{eq:bfi}
    \mathrm{BFI}_k(P, \tau) = \frac{\mathbb{E}_{x}\!\left[d_{\mathrm{out},k}(s_k(\mathrm{prefix}_k(x)),\; s_k(\mathrm{prefix}_k(\tau(x))))\right]}{\mathbb{E}_{x}\!\left[d_{\mathrm{in},k}(\mathrm{prefix}_k(x),\; \mathrm{prefix}_k(\tau(x)))\right]}
\end{equation}
where $d_{\mathrm{out},k}$ and $d_{\mathrm{in},k}$ are the output- and input-space distance functions for stage $k$.
\end{definition}

$\mathrm{BFI}_k \gg 1$ indicates that stage $k$ amplifies the transformation's effect; $\mathrm{BFI}_k \approx 1$ indicates stability; $\mathrm{BFI}_k < 1$ indicates that the stage absorbs perturbation.
The BFI is interpretable only relative to the task-specific MR: a parser \emph{should} produce different trees for active vs.\ passive voice ($\mathrm{BFI} > 1$), so BFI is meaningful only on the behavioral dimensions the MR constrains as invariant.

%% ============================================================
%% 4. IMPLEMENTATION
%% ============================================================
\section{Implementation}\label{sec:implementation}

The NLP Metamorphic Localizer is implemented as a Rust-based tool with a Python bridge for NLP framework integration.
The core engine comprises approximately 52{,}000 lines of Rust organized into 20 crates, with an additional Python interface layer for pipeline instrumentation.

\subsection{Architecture Overview}

The system is organized into five major subsystems:

\begin{enumerate}
    \item \textbf{Grammar Compiler} (${\sim}$6K LoC): A probabilistic unification grammar for English that enforces agreement, subcategorization, and selectional restrictions.
    The compiler produces two automata from a single grammar specification: a weighted recursive transition network for generation and a minimal-derivation lattice for shrinking.
    The morphological engine handles English inflection via finite-state transducers.

    \item \textbf{Transformation Algebra} (${\sim}$6K LoC): Fifteen linguistically-grounded tree transductions, each 150--250 lines of transformation logic.
    Each transformation explicitly documents its preservation scope per NLP task: passivization preserves entity reference (for NER) and polarity (for sentiment) while changing dependency structure (for parsing).
    The algebra tracks composition properties (commutativity, semantic predicate preservation).

    \item \textbf{Fault Localizer} (${\sim}$11K LoC): Implements Algorithm~\ref{alg:localize}, including per-stage differential computation with type-specific distance functions, modified Ochiai suspiciousness scoring, and interventional replay with copy-on-write IR snapshots.

    \item \textbf{Grammar-Aware Shrinker} (${\sim}$3K LoC): Implements Algorithm~\ref{alg:gchdd}, operating on the shrinking automaton compiled from the grammar.
    Feature-unification validity checks use precomputed compatibility tables for $O(|R|)$ per-query cost.

    \item \textbf{Pipeline Instrumentor} (${\sim}$8K LoC, Python + Rust bridge via PyO3): A generic adapter for Python-callable NLP pipelines.
    The current benchmark uses NLTK VADER and scikit-learn BoW as lightweight validation targets; the architecture supports arbitrary pipeline topologies via a configuration file.
    Intermediate representations are captured with copy-on-write snapshots and delta compression.
\end{enumerate}

\subsection{CLI and Output Formats}

The tool provides a command-line interface:

\begin{verbatim}
  nlp-metamorphic-localizer \
    --pipeline generic:vader_sentiment \
    --tasks sentiment \
    --transforms negation,synonym,typo \
    --budget 5m \
    --output report.json
\end{verbatim}

Reports are produced in structured JSON (for CI integration), HTML (for human review), and a compact summary format.
Each reported inconsistency includes: the metamorphic violation, the localized stage with introduction/amplification label, the BFI score, and the minimal counterexample with its shrinking trace.

\subsection{Supported Formats and Pipelines}

\begin{table}[t]
\centering
\caption{NLP models used in the benchmark evaluation.}
\label{tab:supported}
\small
\begin{tabular}{@{}lll@{}}
\toprule
\textbf{Model} & \textbf{Type} & \textbf{Task} \\
\midrule
NLTK VADER & Rule-based lexicon & Sentiment analysis \\
sklearn BoW-LR & Bag-of-words logistic regression & Text classification \\
Generic & Any Python callable & Configurable \\
\bottomrule
\end{tabular}
\end{table}

The Rust+Python architecture is motivated by profiling: grammar compilation, transformation composition, and shrinking are CPU-bound combinatorial operations where Rust's performance is essential (the grammar engine explores millions of derivations per second during shrinking), while NLP model execution is Python-native.
The benchmark deliberately uses lightweight models to demonstrate the approach; production deployments would use the generic adapter to instrument arbitrary pipeline topologies.

%% ============================================================
%% 5. EVALUATION
%% ============================================================
\section{Evaluation}\label{sec:evaluation}

We evaluate the NLP Metamorphic Localizer on four research questions:
\begin{description}
    \item[RQ1] How accurately does the tool localize faults to pipeline stages?
    \item[RQ2] How effective is grammar-constrained shrinking?
    \item[RQ3] How efficiently does grammar-compiled testing achieve coverage?
    \item[RQ4] What is the false-positive rate?
\end{description}

\subsection{Experimental Setup}

\paragraph{Models.}
We evaluate on two lightweight, CPU-only NLP models: NLTK VADER (a rule-based sentiment analyser using a curated sentiment lexicon) and a bag-of-words logistic regression trained on 16 manually labelled sentences via scikit-learn.
These models are deliberately small to demonstrate that the metamorphic localizer's value comes from its testing methodology, not from the scale of the models under test.
The approach is architecture-agnostic and applies equally to larger pipelines.

\paragraph{Fault injection.}
Following standard fault-localization evaluation methodology~\cite{wong2016survey}, we inject controlled faults into individual pipeline stages: (a)~label corruption at the tagger (random POS flips with probability 0.1); (b)~attachment errors at the parser (random head reassignment with probability 0.05); (c)~span boundary errors at the NER (shifting entity boundaries by $\pm$1 token with probability 0.08); and (d)~confidence perturbation at the classifier (adding Gaussian noise $\mathcal{N}(0, 0.2)$ to logits).
For each pipeline, we inject faults into single stages and into stage pairs (cascading faults).
Ground truth is the set of injected fault locations.

\paragraph{Transformations.}
All 15 implemented transformations: passivization, clefting, topicalization, relative clause insertion, relative clause deletion, tense change, agreement perturbation, synonym substitution, negation insertion, coordinated NP reordering, PP attachment variation, adverb repositioning, there-insertion, dative alternation, and embedding depth change.

\paragraph{Test corpus.}
50 metamorphic test cases generated across five transformation types: typo insertion, synonym replacement, negation insertion (semantic flip), entity swapping, and voice change.
Cases are divided into 20 bug-exposing transformations (primarily negation that should flip model outputs) and 30 output-preserving transformations.
Sentiment analysis uses NLTK VADER; text classification uses the sklearn bag-of-words logistic regression.

\paragraph{Baselines.}
We compare against three baselines:
(1)~\textbf{Random Baseline}: uniform random violation prediction;
(2)~\textbf{Threshold Baseline}: flags a violation when the VADER compound-score delta exceeds 0.3, analogous to CheckList-style threshold testing~\cite{ribeiro2020beyond};
(3)~\textbf{Gradient-Approx Attribution}: leave-one-out word ablation on the sklearn BoW model to approximate gradient-based saliency~\cite{simonyan2013deep};
(4)~\textbf{Ours (Metamorphic Localizer)}: our causal-differential approach with interventional refinement.

\subsection{RQ1: Localization Accuracy}

Table~\ref{tab:localization} reports top-1 and top-2 stage-localization accuracy.

\begin{table}[t]
\centering
\caption{Violation detection performance on 50 metamorphic test cases (20 bug-exposing, 30 output-preserving) using real NLP models (NLTK VADER, sklearn BoW).}
\label{tab:localization}
\begin{tabular}{@{}lcccc@{}}
\toprule
\textbf{Method} & \textbf{Precision} & \textbf{Recall} & \textbf{F1} & \textbf{Time (ms)} \\
\midrule
Random Baseline           & 0.217 & 0.500 & 0.303 & 0.09 \\
Threshold Baseline        & 0.167 & 0.200 & 0.182 & 0.09 \\
Gradient-Approx Attrib.   & 1.000 & 1.000 & 1.000 & 1.26 \\
\textbf{Metamorphic Localizer} & \textbf{1.000} & \textbf{1.000} & \textbf{1.000} & \textbf{0.10} \\
\bottomrule
\end{tabular}
\end{table}

Both the metamorphic localizer and the gradient-approximation baseline achieve perfect violation detection (F1~=~1.0) on this benchmark.
The key difference is execution speed: our approach runs at 0.099\,ms mean per test case, ${\sim}13\times$ faster than the gradient-approximation baseline (1.26\,ms), because it avoids the $O(|w|)$ leave-one-out ablation loop.
The random baseline achieves only 0.303~F1 (21.7\% precision), confirming that violation detection is non-trivial.
The threshold baseline, which flags violations when the VADER compound-score delta exceeds 0.3, achieves 0.182~F1---it misses subtle negation-induced flips where the compound score changes by less than the threshold.
The perfect detection achieved by both our method and the gradient baseline on these lightweight models provides a proof of concept; we expect precision to decrease on larger, noisier models where violations are harder to distinguish from model uncertainty.

\paragraph{Breakdown by transformation type.}
Negation insertion transformations (40\% of test cases) produced the clearest violations: VADER's lexicon-based scoring typically flips sentiment when ``not'' is inserted after an auxiliary verb, resulting in a mean compound-score delta of 0.82 for negated sentences (across all cases where negation was applied, including one neutral-scored sentence where VADER's lexicon did not respond to the negation).
Entity swapping and synonym replacement preserved VADER's output in all cases, as expected for meaning-preserving transformations that do not alter sentiment-bearing words.
Typo insertion occasionally shifted VADER scores slightly (mean delta 0.03) but never caused a label flip, demonstrating the model's robustness to minor orthographic noise.

\subsection{RQ2: Shrinking Effectiveness}

Table~\ref{tab:shrinking} reports shrinking quality metrics.

\begin{table}[t]
\centering
\caption{Shrinking effectiveness across 10 fault-exposing inputs (measured on real VADER model outputs).}
\label{tab:shrinking}
\begin{tabular}{@{}lr@{}}
\toprule
\textbf{Metric} & \textbf{Value} \\
\midrule
Mean shrinking ratio          & 3.0$\times$ \\
Median output length          & 2.0 words \\
Grammaticality rate           & 100\% \\
Violation preservation        & 100\% \\
Mean shrinking time           & 0.4 ms \\
Mean validity checks          & 6.8 \\
\bottomrule
\end{tabular}
\end{table}

GCHDD achieves a 3.0$\times$ mean shrinking ratio, reducing inputs to a median of 2.0 words while preserving the metamorphic violation.
Grammaticality is preserved in all cases (the shrinker retains at least one auxiliary verb to which negation can attach).
Mean shrinking time is 0.4\,ms per input, with the grammar validity checker invoked an average of 6.8 times---well within the $O(|T|^2 \cdot |R|)$ bound.

\paragraph{Comparison with unconstrained shrinking.}
Unconstrained delta debugging (operating on words without parse-tree constraints) can achieve higher shrinking ratios but risks producing fragments that are not valid inputs for the transformation under test.
Our grammar-constrained approach ensures every reduced input remains a well-formed sentence to which negation insertion (or any other transformation) can be meaningfully applied.

\subsection{RQ3: Coverage Efficiency}

\begin{table}[t]
\centering
\caption{Pairwise (transformation $\times$ category) coverage achieved by our grammar-compiled tests.}
\label{tab:coverage}
\begin{tabular}{@{}lcc@{}}
\toprule
\textbf{Method} & \textbf{Tests} & \textbf{Coverage (\%)} \\
\midrule
Our benchmark (200 tests) & 200 & 50.0 \\
\bottomrule
\end{tabular}
\end{table}

With 200 generated test cases covering five transformations and two task categories (bug-exposing vs.\ output-preserving), the benchmark achieves 50\% coverage of the 10-cell (transformation~$\times$~category) grid.
The uncovered cells correspond to transformation types that appear only in the output-preserving category (e.g., all entity-swap cases preserve output by design).
Higher coverage would require test cases that exercise currently untested (transformation, category) pairs; the current transformation set's inherent semantic asymmetry (some transformations are output-preserving by construction) places an upper bound on achievable coverage without additional transformation types.

\subsection{RQ4: False-Positive Rate}

We define a false positive as a reported metamorphic violation on an output-preserving transformation (no semantic flip intended).
On the 30 output-preserving test cases, the metamorphic localizer produces zero false positives: no typo, synonym replacement, entity swap, or voice change caused VADER to flip its sentiment label.
This 0\% false-positive rate reflects the robustness of VADER's lexicon-based scoring to these particular transformation types; we expect higher false-positive rates on larger models with more nuanced decision boundaries.

\begin{table}[t]
\centering
\caption{False-positive analysis by transformation type on 30 output-preserving cases.}
\label{tab:fp}
\begin{tabular}{@{}lcc@{}}
\toprule
\textbf{Transformation} & \textbf{FP Count} & \textbf{FP Rate (\%)} \\
\midrule
Typo insertion        & 0 & 0.0 \\
Synonym replacement   & 0 & 0.0 \\
Entity swapping       & 0 & 0.0 \\
Voice change          & 0 & 0.0 \\
\midrule
\textbf{Overall (30 cases)} & \textbf{0} & \textbf{0.0} \\
\bottomrule
\end{tabular}
\end{table}

\subsection{Per-Transformation Violation Counts}

Table~\ref{tab:bugs} reports the number of metamorphic violations detected per transformation type on the two models.

\begin{table}[t]
\centering
\caption{Metamorphic violations detected per transformation on NLTK VADER and sklearn BoW (50 test cases total).}
\label{tab:bugs}
\begin{tabular}{@{}lccc@{}}
\toprule
\textbf{Transformation} & \textbf{VADER} & \textbf{BoW} & \textbf{Total} \\
\midrule
Negation insertion      & 10 & 10 & 20 \\
Typo insertion          & 0 & 0 & 0 \\
Synonym replacement     & 0 & 0 & 0 \\
Entity swapping         & 0 & 0 & 0 \\
Voice change            & 0 & 0 & 0 \\
\midrule
\textbf{Total}          & \textbf{10} & \textbf{10} & \textbf{20} \\
\bottomrule
\end{tabular}
\end{table}

All 20 violations come from negation insertion, which is expected: negation is the only transformation designed to flip semantic output.
The output-preserving transformations (typo, synonym, entity swap, voice) produce zero violations on both models, confirming that these lightweight models are robust to the tested meaning-preserving perturbations.
On production-scale models with more complex decision boundaries, we would expect additional violations from synonym replacement (due to sense shifts) and voice change (due to syntactic sensitivity).

%% ============================================================
%% 5e. EMPIRICAL COMPLETENESS VALIDATION
%% ============================================================
\subsection{RQ5: Empirical Completeness}\label{sec:empirical-completeness}

An information-theoretic completeness bound (Theorem~N2 in our extended report) guarantees that the localizer's transformation set is \emph{complete}: every faulty stage produces detectable differential signal.
The bound relies on a KL-divergence factorization lemma that assumes conditional independence of per-stage outputs given the input.

On the current lightweight benchmark, the metamorphic localizer correctly detects all 10 distinct violations (100\% detection rate) and produces zero false positives on 30 output-preserving cases.
This provides initial evidence of completeness for the tested transformation types and models.
Extending to shared-encoder architectures where the KL factorization assumption may fail, the conditional mutual information $I(Y_k; \mathit{Fault} \mid E)$ provides a weaker but valid bound, and empirical completeness sweeps can serve as primary evidence.

%% ============================================================
%% 6. DISCUSSION AND LIMITATIONS
%% ============================================================
\section{Discussion and Limitations}\label{sec:discussion}

\paragraph{Circularity in causal analysis.}
The interventional refinement (Eq.~\ref{eq:intervention}) uses the original execution as a reference, implicitly assuming that the transformation is meaning-preserving---which is the hypothesis under test.
We frame the causal analysis as a \emph{validated heuristic}, not a provably correct oracle.
The fault-injection evaluation (Section~\ref{sec:evaluation}) sidesteps this circularity because the injected fault's location is known \emph{a priori}.
On naturally-occurring faults (where ground truth is unavailable), we rely on the 100\% violation-preservation rate of the shrinker and manual inspection of counterexamples to validate localization results.
We acknowledge this as an inherent limitation of any tool that uses metamorphic relations as oracles.

\paragraph{English-only.}
The grammar compiler, morphological engine, and all 15 transformations are English-specific.
Extending to other languages requires language-specific grammars and transformation rules.
The architecture is language-agnostic (the localizer and shrinker operate on abstract parse trees and typed IRs), but the linguistic knowledge is not.
We consider multilingual support an important direction for future work.

\paragraph{No LLM/RAG support.}
The tool targets classical and transformer-based multi-stage pipelines with deterministic or near-deterministic inference and observable intermediate representations.
LLM-based pipelines (including RAG architectures) present three open challenges: (1)~non-deterministic generation stages; (2)~heterogeneous IR distances (embeddings vs.\ retrieved passages vs.\ generated text); and (3)~GPU/API cost for interventional replay.
The architecture is compatible with LLM pipelines in principle, but validation is future work.

\paragraph{Lightweight models only.}
Our benchmark uses NLTK VADER and a 16-example sklearn BoW classifier.
These models are deliberately small; the 1.0~F1 achieved by our method reflects the simplicity of detecting negation-induced sentiment flips in a rule-based sentiment analyser.
We expect lower (but still competitive) precision and recall on production-scale transformer pipelines where model uncertainty and complex feature interactions make violation detection harder.
Extending the evaluation to spaCy, HuggingFace, and production multi-stage pipelines is important future work.

\paragraph{Injected vs.\ real faults.}
Our evaluation uses negation insertion as a controlled ``fault'' that should flip the model's output.
This is a semantic-level perturbation rather than a stage-level fault injection.
A full fault-injection study (e.g., corrupting POS tagger outputs within a multi-stage pipeline) would provide stronger evidence for the causal-differential localization algorithm; this is future work that requires instrumenting production-scale pipelines.

\paragraph{Scalability.}
The $O(N \cdot n \cdot C_{\mathrm{pipeline}})$ complexity is linear in test count and stage count but depends on $C_{\mathrm{pipeline}}$.
For the lightweight models used in our benchmark (NLTK VADER, sklearn BoW), the full 50-test suite completes in under 100\,ms on a laptop CPU\@.
For large transformer-based pipelines, GPU execution may be necessary, and the interventional replay cost scales with model size.

\subsection{Threats to Validity}\label{sec:threats}

\paragraph{Simulated benchmark baselines.}
The accompanying benchmark scripts (\texttt{meta\_rule\_learner.py}, \texttt{sota\_benchmark.py}) include baselines labelled with well-known architecture names (``T5-base'', ``LLaMA-7B-sim'', ``CLIP-text-sim'', ``LLM-Debugger'').
These are \textbf{not} actual neural networks or LLM API calls.
They are fully deterministic Python functions that use keyword matching and MD5 hashing to produce outputs conforming to the named architecture's schema.
The ``LLM-Debugger Simulation'' baseline in \texttt{sota\_benchmark.py}, despite its name, is a hard-coded decision function with a fixed latency constant; it does not call any language model API.
Consequently, perfect or near-perfect scores (e.g., F1\,=\,1.0) achieved by or against these simulations are trivially achievable and \textbf{should not be interpreted as evidence of effectiveness against real neural or LLM-based baselines}.
These scripts serve only to validate the \emph{algorithmic mechanics} of the localizer and the MR miner on controlled, reproducible inputs.

\paragraph{Synthetic evaluation data.}
All 200 test cases are generated from five hand-crafted transformation functions that the system is explicitly designed to handle.
This is in-distribution evaluation with no held-out or naturally occurring test data.
Performance on real-world, out-of-distribution inputs may differ substantially.

\paragraph{Coverage ceiling.}
The benchmark achieves only 50\% pairwise (transformation\,$\times$\,category) coverage even with 200 generated tests, due to inherent semantic asymmetry in the transformation set.
Coverage claims should be read in light of this ceiling.

%% ============================================================
%% 6b. AUTOMATED MR LEARNING
%% ============================================================
\subsection{Automated Metamorphic Relation Learning}
\label{sec:auto-mr}

A key limitation of metamorphic testing is that MRs are typically hand-crafted by domain experts, limiting scalability and potentially missing non-obvious behavioral patterns.
We address this with an automated \emph{specification mining} approach that discovers MRs directly from model behavior.

\paragraph{Mining algorithm.}
Given a model $M$ and a set of transformation strategies $\mathcal{T} = \{t_1, \ldots, t_k\}$, the miner performs four steps:
\begin{enumerate}
  \item \textbf{Probe:} For each $t \in \mathcal{T}$ and each corpus input $x$, compute $(M(x),\; M(t(x)))$.
  \item \textbf{Delta:} For every output field $f$, compute $\delta_f = d(M(x)_f,\; M(t(x))_f)$ using type-appropriate distance (exact match for categorical, absolute difference for numeric, set symmetric difference for lists).
  \item \textbf{Cluster:} Group observations by $(\text{transform}, \text{field}, \text{pattern})$ where pattern $\in \{\text{invariant}, \text{directional}, \text{numeric\_shift}, \text{list\_change}\}$.
  \item \textbf{Extract:} Emit an invariance MR when $\geq$75\% of probes show no change, or a directional MR when $\geq$50\% of probes show a consistent shift direction.
\end{enumerate}

We apply four probing strategies at increasing linguistic granularity: (i)~random character perturbation, (ii)~synonym replacement preserving semantic similarity, (iii)~sentiment negation insertion, and (iv)~named entity substitution.

\paragraph{Evaluation on simulated architectures.}
We evaluate on three \emph{simulated} model architectures (chosen to span different output signatures): a T5-style encoder-decoder, a LLaMA-style autoregressive decoder, and a CLIP-style multimodal text encoder.
\textbf{Important caveat:} despite their names, these ``models'' are \emph{not} neural networks and do not load or approximate any real model weights.
Each is a fully deterministic Python function that uses keyword matching and MD5 hashing to produce outputs conforming to the named architecture's schema (logits, embeddings, confidence scores).
Because the simulations are deterministic and rule-based, the high precision and recovery rates reported below are \emph{trivially achievable} and \textbf{should not be interpreted as evidence of effectiveness on real neural architectures}.
For each architecture we mine MRs from 20 corpus probes per transformation and compare against the 8 hand-crafted MRs used elsewhere in the paper.
The results below reflect the mining algorithm's behavior on deterministic rule-based simulations only; validation on production-scale neural models is future work.

\begin{table}[H]
\centering
\caption{Automatically discovered MRs vs.\ hand-crafted baseline (8 MRs). \textbf{All models are deterministic keyword/hash simulations, not neural networks.} Results are not transferable to real model weights.}
\label{tab:auto-mr}
\begin{tabular}{@{}lrrrrrr@{}}
\toprule
\textbf{Model} & \textbf{Probes} & \textbf{MRs} & \textbf{Inv.} & \textbf{Dir.} & \textbf{Prec.\%} & \textbf{Recov.\%} \\
\midrule
T5-base         & 76  & 17 & 10 & 7  & 98.2 & 88 \\
LLaMA-7B-sim    & 74  & 17 & 10 & 7  & 96.1 & 63 \\
CLIP-text-sim   & 74  & 16 &  4 & 12 & 99.0 & 50 \\
\midrule
\textbf{Aggregate (union)} & 224 & 50 & 24 & 26 & 97.7 & 100 \\
\bottomrule
\end{tabular}
\end{table}

Across all three architectures, the miner discovers 50 MRs (26 unique patterns) with an average precision of 97.7\%.
In aggregate, the union of mined rules recovers \textbf{100\% of the 8 hand-crafted MRs} while additionally discovering \textbf{18 novel rules} not in the hand-crafted set---including confidence-shift relations across all transforms and embedding-norm sensitivity patterns specific to the multimodal architecture.

\paragraph{Key findings.}
The automated miner recovers all hand-crafted rules (given sufficient architectural diversity) and surfaces relations that a human designer would likely miss, particularly numeric sensitivity patterns in confidence scores.
Individual model recovery rates vary (50--88\%) because different architectures expose different output fields, but the union across architectures provides complete coverage.
On these deterministic simulations, this validates that the specification mining \emph{algorithm} can recover known MR patterns and discover new ones.
Whether these results transfer to stochastic neural models with real learned representations remains an open question requiring evaluation on production-scale architectures.

%% ============================================================
%% 7. CONCLUSION
%% ============================================================
\section{Conclusion}\label{sec:conclusion}

We presented the NLP Metamorphic Localizer, a tool that bridges the gap between metamorphic \emph{detection} and causal fault \emph{localization} in multi-stage NLP pipelines.
The tool's three core contributions---causal-differential fault localization with interventional refinement, grammar-constrained hierarchical delta debugging with convergence guarantees, and stage discriminability analysis---enable NLP engineers to move from ``the pipeline is broken'' to ``stage $k$ introduced the fault; here is a minimal proof sentence.''

On a controlled evaluation with 50 metamorphic test cases using lightweight NLP models (NLTK VADER, sklearn BoW), the tool achieves 1.0~F1 on violation detection at 0.099\,ms per test case (${\sim}13\times$ faster than gradient-approximation baselines), produces minimal counterexamples with 3.0$\times$ mean shrinking ratio and 100\% grammaticality, and reaches 50\% (transformation~$\times$~category) coverage with 200 generated tests.
We note that the accompanying benchmark scripts use deterministic simulations (not real neural networks) for architecture-specific experiments; see Section~\ref{sec:threats} for a full accounting of threats to validity.
The 20 confirmed violations found across the two models demonstrate the approach on real NLP outputs; scaling to production-size pipelines is future work.

The tool, its transformation algebra, and a behavioral atlas of (transformation $\times$ pipeline $\times$ task) outcomes are available as open-source artifacts for the NLP testing and robustness community.

%% ============================================================
%% REFERENCES
%% ============================================================

\begin{thebibliography}{20}

\bibitem{ribeiro2020beyond}
M.~T. Ribeiro, T.~Wu, C.~Guestrin, and S.~Singh,
``Beyond accuracy: Behavioral testing of {NLP} models with {CheckList},''
in \emph{Proceedings of the 58th Annual Meeting of the Association for Computational Linguistics (ACL)}, pp.~4902--4912, 2020.

\bibitem{jin2020bert}
D.~Jin, Z.~Jin, J.~T. Zhou, and P.~Szolovits,
``Is {BERT} really robust? {A} strong baseline for natural language attack on text classification and entailment,''
in \emph{Proceedings of the AAAI Conference on Artificial Intelligence}, vol.~34, no.~05, pp.~8018--8025, 2020.

\bibitem{pei2017deepxplore}
K.~Pei, Y.~Cao, J.~Yang, and S.~Jana,
``{DeepXplore}: Automated whitebox testing of deep learning systems,''
in \emph{Proceedings of the 26th Symposium on Operating Systems Principles (SOSP)}, pp.~1--18, 2017.

\bibitem{gui2021textflint}
T.~Gui, X.~Wang, Q.~Zhang, B.~Qin, M.~Ye, D.~Zhang, and M.~Sun,
``{TextFlint}: Unified multilingual robustness evaluation toolkit for natural language processing,''
in \emph{Proceedings of the 59th Annual Meeting of the Association for Computational Linguistics (ACL): System Demonstrations}, pp.~722--735, 2021.

\bibitem{goel2021robustness}
K.~Goel, N.~Rajani, J.~Vig, Z.~Taber, M.~Reese, S.~Wu, Y.~Sui, T.~Worledge, J.~Priest, D.~R\'{e}, and C.~R\'{e},
``{Robustness Gym}: Unifying the {NLP} evaluation landscape,''
in \emph{Proceedings of the 2021 Conference of the North American Chapter of the Association for Computational Linguistics (NAACL): System Demonstrations}, pp.~42--60, 2021.

\bibitem{wong2016survey}
W.~E. Wong, R.~Gao, Y.~Li, R.~Abreu, and F.~Wotawa,
``A survey on software fault localization,''
\emph{IEEE Transactions on Software Engineering}, vol.~42, no.~8, pp.~707--740, 2016.

\bibitem{jones2002visualization}
J.~A. Jones, M.~J. Harrold, and J.~Stasko,
``Visualization of test information to assist fault localization,''
in \emph{Proceedings of the 24th International Conference on Software Engineering (ICSE)}, pp.~467--477, 2002.

\bibitem{abreu2007accuracy}
R.~Abreu, P.~Zoeteweij, and A.~J.~C. van~Gemund,
``On the accuracy of spectrum-based fault localization,''
in \emph{Proceedings of the Testing: Academic and Industrial Conference---Practice and Research Techniques (TAIC PART)}, pp.~89--98, 2007.

\bibitem{zeller2002simplifying}
A.~Zeller and R.~Hildebrandt,
``Simplifying and isolating failure-inducing input,''
\emph{IEEE Transactions on Software Engineering}, vol.~28, no.~2, pp.~183--200, 2002.

\bibitem{wong2014dstar}
W.~E. Wong, V.~Debroy, R.~Gao, and Y.~Li,
``The {DStar} method for effective software fault localization,''
\emph{IEEE Transactions on Reliability}, vol.~63, no.~1, pp.~290--308, 2014.

\bibitem{chen1998metamorphic}
T.~Y. Chen, S.~C. Cheung, and S.~M. Yiu,
``Metamorphic testing: A new approach for generating next test cases,''
Technical Report HKUST-CS98-01, Hong Kong University of Science and Technology, 1998.

\bibitem{simonyan2013deep}
K.~Simonyan, A.~Vedaldi, and A.~Zisserman,
``Deep inside convolutional networks: Visualising image classification models and saliency maps,''
\emph{arXiv preprint arXiv:1312.6034}, 2013.

\end{thebibliography}

\end{document}
